\documentclass[11pt]{article}

\usepackage[a4paper,margin=1in]{geometry}
\usepackage{amsmath,amssymb,amsthm,mathtools}
\usepackage{enumitem}
\usepackage[hidelinks]{hyperref}

\newtheorem{theorem}{Theorem}
\newtheorem{proposition}[theorem]{Proposition}
\newtheorem{lemma}[theorem]{Lemma}
\newtheorem{corollary}[theorem]{Corollary}
\theoremstyle{definition}
\newtheorem{definition}[theorem]{Definition}
\newtheorem{remark}[theorem]{Remark}
\newtheorem{problem}[theorem]{Problem}

\newcommand{\R}{\mathbb R}
\newcommand{\Om}{\Omega}
\newcommand{\HH}{\mathcal H}
\newcommand{\eps}{\varepsilon}
\newcommand{\dT}{\delta_T}
\newcommand{\AF}{\mathcal A_F}
\newcommand{\Per}{\mathrm{Per}}
\newcommand{\esssup}{\operatorname*{ess\,sup}}
\newcommand{\norm}[1]{\left\|#1\right\|}
\newcommand{\abs}[1]{\left|#1\right|}

\title{From Two-Sided Density to a Radial Graph:\\
       what the realization step really needs (and what it does not)}
\author{}
\date{}

\begin{document}
\maketitle

\begin{abstract}
We isolate the two independent inputs of the ``realization'' step in the
selection-free route to quantitative Saint--Venant: (A) an
\emph{inclusion} $B_{r_-}(z)\subset\Om\subset B_{r_+}(z)$, and (B) the
\emph{realization} of $\Om$ as a radial graph with a \emph{quadratic}
pocket estimate $|G\setminus\Om|\lesssim(P(\Om)-2\pi)^2$. The first
version of this note claimed that (B) holds for any set of finite
perimeter. \textbf{That claim is false}, and we begin by giving the
explicit counterexample (a ``lollipop'': a wide cap overhanging a thin
stem), which shows that finite perimeter alone never controls the pocket
volume quadratically --- not even inside an arbitrarily thin annulus, and
not even under two-sided density. The correct hypothesis is a purely
measure-theoretic \emph{no-overhang} condition,
$\HH^1(\partial^*G\setminus\partial^*E)=0$, which is far weaker than
$C^{1,\gamma}$ but strictly stronger than finite perimeter. Under
no-overhang we prove the perimeter decomposition and the sharp quadratic
pocket bound (Theorem~\ref{thm:realize}); this is the honest GMT form of
the Plan~4 pocket lemma, and it does retire the $C^{1,\gamma}$ hypothesis
\emph{provided no-overhang is available}. We then analyse precisely what
input (A) provides: an annulus gives the trivial linear bound
$|G\setminus\Om|\lesssim\eta$ unconditionally, while two-sided density
plus tentacle exclusion removes overhangs whose \emph{necks} sit above
scale $r_*$, leaving a genuine residual of sub-$r_*$ overhangs. The
upshot is a clean trisection of the programme into inclusion, no-overhang,
and the radial-graph theorem, with the residual identified honestly.
\end{abstract}

\section{The two inputs of realization}

Fix the planar normalisation $|\Om|=\pi$, comparison ball $B=B_1$. The
goal of the realization step is to replace a (possibly rough) competitor
$\Om$ by a \emph{radial graph} $G$ with $|G\,\triangle\,\Om|$ controlled
and $\dT(G)$ controlled, so that the small-amplitude radial-graph theorem
$(G)$ applies. Two logically separate things are needed:

\begin{enumerate}[label=\textup{(\Alph*)}]
\item\label{inputA} \emph{Inclusion.} A centre $z$ and radii
$r_-<r_+$ with
$B_{r_-}(z)\subset\Om\subset B_{r_+}(z)$.
\item\label{inputB} \emph{Realization.} From \ref{inputA}, the radial
hull $G$ (star-shaped envelope from $z$) is a radial graph, and
$|G\setminus\Om|$ is controlled \emph{quadratically} by the perimeter
deficit.
\end{enumerate}

The historical role of selection was to make $\partial\Om$ regular enough
($C^{1,\gamma}$) that both \ref{inputA} (via BNST) and \ref{inputB} (via
the polar-slicing pocket lemma, which used Sard and the implicit function
theorem) could run. The question this note addresses is: \emph{how much
regularity does \ref{inputB} actually need?} The answer is more subtle
than ``none'': it needs a no-overhang condition. We first show, by an
explicit example, why finite perimeter is not enough; then we prove the
sharp statement under no-overhang; then we examine what (A) supplies.

\section{Input (A): inclusion from density, recalled}
\label{sec:inclusion}

We recall the mechanism of \texttt{almost\_serrin\_discussion.tex}. Let
$u$ be the torsion function of $\Om$, $c_\Om=|\Om|/P(\Om)$, and define the
Serrin and isoperimetric deficits
\[
   \eps^2:=\int_{\partial\Om}\bigl(|\nabla u|-c_\Om\bigr)^2\,d\HH^{1},
   \qquad
   \delta_{\mathrm{iso}}:=\frac{P(\Om)}{2\sqrt{\pi|\Om|}}-1 .
\]
The local flux identity (integrate $-\Delta u=1$ over $\Om\cap B_r(x)$)
combined with Cauchy--Schwarz against $\eps$ and the relative
isoperimetric inequality in balls yields the \emph{tentacle-exclusion}
estimate: for $x\in\partial\Om$,
\[
   \eps_r^2\ \gtrsim\ \abs{\Om\cap B_r(x)}^{1/2}
   \quad\text{when }\abs{\Om\cap B_r(x)}\ll r^2,
\]
so small Serrin deficit forbids thin interior tentacles above scale
$r_*\sim\eps^{2}$; quantitative isoperimetry forbids exterior fjords above
scale $\delta_{\mathrm{iso}}^{1/4}$. The outcome is the two-sided density
bound
\begin{equation}\label{eq:density}
   c_0\,r^2\ \le\ \abs{\Om\cap B_r(x)}\ \le\ (1-c_0)\,r^2,
   \qquad x\in\partial\Om,\ \ r\in[r_*,r_0],
\end{equation}
with $r_*\sim\eps^{2}+\delta_{\mathrm{iso}}^{1/4}$, and (via an adapted
Magnanini--Poggesi step) the inclusion
\begin{equation}\label{eq:annulus}
   B_{\rho_i}(z)\subset\Om\subset B_{\rho_e}(z),
   \qquad
   \rho_e-\rho_i\ \le\ \eta:=C\bigl(\eps^{\theta}+r_*\bigr).
\end{equation}
We treat \eqref{eq:density}--\eqref{eq:annulus} as the \emph{output} of
input (A); its proof is the content of
\texttt{almost\_serrin\_discussion.tex} (with the open points there). What
matters below is only that (A) delivers an inclusion with some width
$\eta$, a finite perimeter $P(\Om)<\infty$ (immediate from
\eqref{eq:density}, Ahlfors regularity), and the \emph{scale} $r_*$ above
which density holds.

\section{Why finite perimeter is not enough}
\label{sec:counterexample}

We work in $\R^2$ with polar coordinates about $0$. For a set $E$ of
finite perimeter with $B_{r_-}(0)\subset E\subset B_{r_+}(0)$, define the
radial slices, the profile, and the radial hull
\begin{equation}\label{eq:hull}
   E_\theta:=\{r:re^{i\theta}\in E^{(1)}\},\qquad
   R(\theta):=\esssup E_\theta,\qquad
   G:=\{re^{i\theta}:0\le r<R(\theta)\},
\end{equation}
and the pocket set $Q:=G\setminus E$, with connected components
$\{Q_j\}_{j\in J}$. The realization step rests on the would-be
\emph{perimeter decomposition}
\begin{equation}\label{eq:wrong}
   P(G)+\sum_{j\in J}P(Q_j)\ \le\ P(E)
   \qquad(\text{\emph{false} in general}).
\end{equation}
We exhibit a finite-perimeter $E$ for which \eqref{eq:wrong} fails by a
fixed positive amount.

\begin{proposition}[Finite perimeter is insufficient]\label{prop:lollipop}
Let $0<a<b<R^+$ and $0<w'<w<2\pi$, and set (the ``lollipop'')
\[
   E=B_a\ \cup\ \underbrace{\{a\le r\le b,\ |\theta|\le w'/2\}}_{\text{stem }S}
   \ \cup\ \underbrace{\{b\le r\le R^+,\ |\theta|\le w/2\}}_{\text{cap }C}.
\]
Then $E$ has finite perimeter, $B_a\subset E\subset B_{R^+}$, its radial
hull is $G=\{|\theta|<w/2,\ r<R^+\}\cup\{|\theta|>w/2,\ r<a\}$ with two
pockets $Q^\pm=\{w'/2<\pm\theta<w/2,\ a<r<b\}$, and
\begin{equation}\label{eq:excess}
   P(G)+P(Q^+)+P(Q^-)\ -\ P(E)\ =\ 4\,(b-a)\ >\ 0 .
\end{equation}
\end{proposition}

\begin{proof}
Split every reduced boundary into its \emph{radial} part (circular arcs,
normal $\nu=\pm e_r$) and its \emph{angular} part (radial walls, normal
$\nu=\pm e_\theta$); all three sets here are polar rectangles, so each
perimeter splits exactly, $P(F)=P_r(F)+P_\theta(F)$.

\emph{Radial part cancels.} Slicing radially, $P_r(F)=\int_{S^1}\sum_{\rho
\in\partial^*F_\theta}\rho\,d\theta$. For a.e.\ $\theta$ the outer crossing
of $G_\theta$ is $R(\theta)$, the crossings of $Q_\theta$ are the inner
gap endpoints of $E_\theta$, and together they are exactly the crossings of
$E_\theta$ with the same radii (the slice balance: $\#\partial^*G_\theta+
\#\partial^*Q_\theta=\#\partial^*E_\theta$, weighted by the radii). Hence
$P_r(G)+P_r(Q^+)+P_r(Q^-)=P_r(E)$.

\emph{Angular part does not.} Directly:
\[
   P_\theta(E)=\underbrace{2(b-a)}_{\text{stem walls }|\theta|=w'/2}
              +\underbrace{2(R^+-b)}_{\text{cap walls }|\theta|=w/2}
              =2(R^+-a),
\]
\[
   P_\theta(G)=\underbrace{2(R^+-a)}_{\text{hull walls }|\theta|=w/2,\ a<r<R^+},
   \qquad
   P_\theta(Q^+)+P_\theta(Q^-)=\underbrace{4(b-a)}_{\text{walls }|\theta|=w'/2\text{ and }|\theta|=w/2,\ a<r<b}.
\]
Subtracting, $P_\theta(G)+P_\theta(Q^+)+P_\theta(Q^-)-P_\theta(E)=
2(R^+-a)+4(b-a)-2(R^+-a)=4(b-a)$, and adding the (vanishing) radial
discrepancy gives \eqref{eq:excess}.
\end{proof}

\begin{remark}[Diagnosis: the overhang]\label{rem:overhang}
The culprit is the cap wall $\{|\theta|=w/2,\ a<r<b\}$, the lateral face
of the pocket $Q^\pm$ sitting \emph{underneath} the cap's overhang
($w>w'$). Both sides of this wall --- the pocket (empty $E$) and the
exterior (empty $E$) --- have $E$-density $0$, so it lies in neither
$\partial^*E$. Yet it is charged once to $P(G)$ (a radial wall of the
hull, where $R$ jumps from $R^+$ to $a$) and once to $P(Q^\pm)$. It is
double-counted on the left of \eqref{eq:wrong} and absent on the right.
This is the exact mechanism by which my earlier ``angular part''
inequality failed.
\end{remark}

\begin{remark}[Thin annuli and density do not help]\label{rem:thin}
Two natural rescues both fail.
\begin{enumerate}[label=\textup{(\roman*)}]
\item \emph{Thin annulus.} Take $a=1$, $b=1+\tfrac\eta2$, $R^+=1+\eta$, so
$E$ lies in the annulus $\{1<r<1+\eta\}$ of width $\eta$, while $w$ stays
order $1$ and $w'\to0$. Then the pocket volume is
$|Q^+|+|Q^-|\approx(w-w')\,a\,(b-a)\sim\eta$, whereas the perimeter
deficit is $P(E)-2\pi\sim\eta$, so $(P(E)-2\pi)^2\sim\eta^2\ll\eta\sim
|G\setminus E|$. The quadratic pocket bound is violated by a factor
$\eta^{-1}$, \emph{inside an arbitrarily thin annulus}.
\item \emph{Two-sided density.} The lollipop is a Lipschitz domain; on
$\partial E$ the density is $\approx\tfrac12$ generically and
$\approx\tfrac14$ at the finitely many corners, so \eqref{eq:density}
holds at every boundary point and every scale (down to the corner scale).
Density does not see the overhang, because the offending wall is not a
boundary point of $E$ at all.
\end{enumerate}
Thus neither the inclusion of (A) nor its density bound is, by itself,
enough for the quadratic realization. Something must forbid overhangs.
\end{remark}

\section{Input (B): realization under no overhang}
\label{sec:realize}

We isolate the exact condition that makes the decomposition true.

\begin{definition}[No overhang]\label{def:noov}
A set of finite perimeter $E$ with $B_{r_-}(0)\subset E\subset B_{r_+}(0)$
has \emph{no overhang about $0$} if its radial hull $G$ from \eqref{eq:hull}
satisfies
\[
   \HH^1\bigl(\partial^*G\setminus\partial^*E\bigr)=0 .
\]
Equivalently: $\HH^1$-a.e.\ point of the hull boundary is a reduced
boundary point of $E$ with the same outer normal --- the outer profile of
$E$ is ``self-supporting,'' $E$ being full immediately beneath each piece
of $\partial^*G$.
\end{definition}

No overhang is implied by, but far weaker than, $C^{1,\gamma}$
star-shapedness; and it is strictly stronger than finite perimeter, the
lollipop being the obstruction (there $\partial^*G$ contains the cap-wall
segments $\{|\theta|=w/2,\ a<r<b\}$, which are not in $\partial^*E$).

\begin{lemma}[Slice structure]\label{lem:slice}
For $\HH^1$-a.e.\ $\theta\in S^1$ the slice $E_\theta$ is a
one-dimensional set of finite perimeter with
$(0,r_-]\subset E_\theta\subset(0,r_+]$, with essential boundary the
finite set $\{r_1<\cdots<r_{2m+1}\}$ ($m=m(\theta)$),
\[
   E_\theta=(0,r_1)\cup(r_2,r_3)\cup\cdots\cup(r_{2m},r_{2m+1}),\qquad
   R(\theta)=r_{2m+1},
\]
$G_\theta=(0,R(\theta))$, and the pocket slice
$Q_\theta=(r_1,r_2)\cup\cdots\cup(r_{2m-1},r_{2m})$. Moreover $R\in
BV(S^1)$ with values in $[r_-,r_+]$, $G$ is of finite perimeter, and $G$
is a radial graph about $0$.
\end{lemma}

\begin{proof}
By Vol'pert's slicing theorem \cite[Thm.~3.108]{AFP}, for $\HH^1$-a.e.\
$\theta$ the slice $E_\theta$ is one-dimensional of finite perimeter and
$\partial^*E_\theta=(\partial^*E)_\theta$; in dimension one a set of finite
perimeter is, up to a null set, a finite union of disjoint open intervals.
The inclusion $B_{r_-}\subset E\subset B_{r_+}$ forces
$(0,r_-]\subset E_\theta\subset(0,r_+]$, so the first interval starts at
$0$ and the outermost crossing is $R(\theta)=r_{2m+1}$. The interval
structure gives $G_\theta$, $Q_\theta$. That $R\in BV(S^1)$ is the
statement that the upper envelope of the $BV$ slice family has bounded
variation; $G$, the polar subgraph of $R\in BV(S^1)$ with values in
$[r_-,r_+]$, is then of finite perimeter and is by construction a radial
graph.
\end{proof}

\begin{lemma}[Perimeter decomposition under no overhang]\label{lem:perim}
If $E$ has no overhang about $0$ (Definition~\ref{def:noov}), with hull
$G$ and pockets $\{Q_j\}_{j\in J}$, then, up to $\HH^1$-null sets,
$\partial^*G$ and the $\partial^*Q_j$ are pairwise disjoint subsets of
$\partial^*E$. Consequently
\begin{equation}\label{eq:perim-dec}
   P(G)+\sum_{j\in J}P(Q_j)\ \le\ P(E).
\end{equation}
\end{lemma}

\begin{proof}
Write $G=E\sqcup Q$ mod $\mathcal L^2$-null, with $Q=\bigsqcup_jQ_j$ and
the $Q_j$ the (essentially disjoint, mutually density-$0$) components of
the pocket set.

\emph{$\partial^*G\subset\partial^*E$.} This is the no-overhang
hypothesis.

\emph{Each $\partial^*Q_j\subset\partial^*E$.} Let $x\in\partial^*Q_j$, so
$Q_j$ has density $\tfrac12$ at $x$. If $x\in\partial^*G$ then
$x\in\partial^*E$ by no-overhang, so assume $x\notin\partial^*G$. Since
$Q_j\subset G$, density $\tfrac12$ of $Q_j$ at $x$ forces $G$-density
$\ge\tfrac12$, hence (as $x\notin\partial^*G$, where $G$ would have density
$\tfrac12$) $G$ has density $1$ at $x$. The other pockets $Q_k$ ($k\ne j$)
have density $0$ at $x$ (mutually density-$0$ components). Therefore
$E=G\setminus Q$ has density $1-\tfrac12-0=\tfrac12$ at $x$, with
$\nu_E(x)=-\nu_{Q_j}(x)$; thus $x\in\partial^*E$.

\emph{Disjointness.} If $x\in\partial^*G\cap\partial^*Q_j$ then $G$ and
$Q_j\subset G$ both have density $\tfrac12$, so $E\cap B$ has density $0$
near $x$: $x$ separates $Q_j$ (inside $G$) from $\R^2\setminus G$ with no
$E$ in between, i.e.\ $x\in\partial^*G\setminus\partial^*E$, a null set by
no-overhang. The $\partial^*Q_j$ are mutually $\HH^1$-disjoint as reduced
boundaries of mutually density-$0$ sets.

\emph{Conclusion.} The sets $\partial^*G,\ \partial^*Q_1,\partial^*Q_2,
\dots$ are, mod $\HH^1$-null, pairwise disjoint subsets of $\partial^*E$,
so
\[
   P(G)+\sum_jP(Q_j)
   =\HH^1\!\Big(\partial^*G\sqcup\bigsqcup_j\partial^*Q_j\Big)
   \le\HH^1(\partial^*E)=P(E). \qedhere
\]
\end{proof}

\begin{theorem}[Radial-graph realization under no overhang]\label{thm:realize}
Let $E\subset\R^2$ have finite perimeter, $|E|=\pi$,
$B_{r_-}(0)\subset E\subset B_{r_+}(0)$, the perimeter bound $P(E)\le
2\pi+\eps$, and \emph{no overhang} about $0$. Let $G$ be the radial hull.
Then:
\begin{enumerate}[label=\textup{(\roman*)}]
\item $G$ is a radial graph about $0$, $G=\{re^{i\theta}:0\le r<R(\theta)\}$
with $R\in BV(S^1)$, $r_-\le R\le r_+$;
\item the pocket volume obeys the sharp quadratic bound
\begin{equation}\label{eq:pocket}
   \abs{G\setminus E}\ \le\ \frac{\eps^2}{4\pi}.
\end{equation}
\end{enumerate}
\end{theorem}

\begin{proof}
(i) is Lemma~\ref{lem:slice}. For (ii) set $v:=|G\setminus E|$,
$h:=\sum_jP(Q_j)$. By planar isoperimetry applied to $G$ (area $\pi+v$),
$P(G)\ge 2\sqrt{\pi(\pi+v)}\ge2\pi$. Combining with
Lemma~\ref{lem:perim} and $P(E)\le2\pi+\eps$,
\[
   2\pi+\eps\ \ge\ P(E)\ \ge\ P(G)+h\ \ge\ 2\pi+h,
\]
so $h\le\eps$. Apply isoperimetry to each pocket, $|Q_j|\le
P(Q_j)^2/(4\pi)$, and super-additivity $\sum a_j^2\le(\sum a_j)^2$ for
$a_j\ge0$:
\[
   v=\sum_j|Q_j|\le\frac1{4\pi}\sum_jP(Q_j)^2
   \le\frac1{4\pi}\Big(\sum_jP(Q_j)\Big)^2=\frac{h^2}{4\pi}
   \le\frac{\eps^2}{4\pi}. \qedhere
\]
\end{proof}

\begin{remark}[What survives, what does not]\label{rem:survives}
The squaring mechanism --- isoperimetry per pocket plus super-additivity ---
never needed regularity, and indeed needs no overhang only through
Lemma~\ref{lem:perim}. So the precise content is: \emph{the only place any
hypothesis beyond finite perimeter enters realization is the perimeter
decomposition, and there the necessary and (essentially) sufficient
condition is no-overhang}. The earlier note conflated ``regularity-free''
with ``hypothesis-free''; the lollipop (Proposition~\ref{prop:lollipop})
shows the two are different. What is genuinely regularity-free is the
\emph{passage from no-overhang to the quadratic bound}; what is not free
is no-overhang itself.
\end{remark}

\subsection*{An unconditional fallback}

Without any no-overhang hypothesis one still has, from the inclusion alone,
the trivial linear bound
\begin{equation}\label{eq:trivial}
   \abs{G\setminus E}\ \le\ \pi(\rho_e^2-\rho_i^2)\ \le\ 2\pi\rho_e\,\eta ,
\end{equation}
since both $E$ and $G$ lie in the annulus $\{\rho_i<r<\rho_e\}$ outside
$B_{\rho_i}$. This is the bound that holds for the lollipop, and it is the
only one available there. The whole point of realization is to upgrade
\eqref{eq:trivial} from $O(\eta)$ to $O(\eta^2)$, and
Theorem~\ref{thm:realize} does so exactly under no-overhang.

\section{Does input (A) supply no overhang?}
\label{sec:Agivesnoov}

This is the crux, and the honest answer is: \emph{partially, above scale
$r_*$}. Two-sided density does not by itself remove overhangs
(Remark~\ref{rem:thin}(ii)); what removes them is the \emph{tentacle
exclusion} that produced density in the first place, applied to the
\emph{neck} of an overhang.

\begin{lemma}[Necks of overhangs are tentacles]\label{lem:neck}
Suppose $\Om$ satisfies the inclusion \eqref{eq:annulus} and the tentacle
exclusion estimate of \S\ref{sec:inclusion} at scale $r_*$. Let $Q$ be a
pocket of the radial hull whose lateral wall is an overhang, i.e.\ a
component of $\partial^*G\setminus\partial^*\Om$, and let $\ell$ be the
angular width of $Q$ at its inner radius. If $\ell\ge r_*$ then the cap
producing the overhang is attached to $\Om\setminus B_{\rho_i}$ through a
neck of width $\ge c\,r_*$; more precisely, an overhang with neck width
$w'<c_1 r_*$ is excluded, because the neck region
$\{a<r<b,\ |\theta-\theta_0|<w'\}$ is an interior tentacle:
$|\Om\cap B_{r}(x)|\le C\,w'\,r\ll r^2$ for $x$ at the neck and
$r\in[w',r_*]$, forcing $\eps_r^2\gtrsim r$ and hence Serrin deficit
$\gtrsim r_*$.
\end{lemma}

\begin{proof}[Proof sketch]
An overhang exists only if the cap (full $E$ at $b<r<R^+$, angular width
$w$) is fed from below through a strictly thinner stem (width $w'<w$): some
neck must connect the cap to the bulk $\Om\setminus B_{\rho_i}$, else the
cap is a separate component and contributes a full isoperimetric deficit.
At a point $x$ on the neck, for scales $w'\le r\le \min(b-a,r_*)$ the slab
$\Om\cap B_r(x)$ is contained in the stem strip of width $w'$, so
$|\Om\cap B_r(x)|\le C w' r\le C(w'/r)\,r^2\ll r^2$ once $w'\ll r$. The
tentacle-exclusion estimate $\eps_r^2\gtrsim|\Om\cap B_r(x)|^{1/2}$ is then
not the sharp form; rather one uses the perimeter lower bound
$P(\Om;B_r(x))\ge c\,r$ (two long stem walls of length $\ge r$) against the
flux identity to get $\eps_r^2\gtrsim r$, so any neck with $w'\ll r$
present at a scale $r\le r_*$ would force $\eps^2\gtrsim r_*$, contradicting
smallness once $r$ is taken $\sim r_*$. Hence necks of width $<c_1r_*$ are
excluded above scale $r_*$.
\end{proof}

\begin{remark}[The residual is genuine and sub-$r_*$]\label{rem:residual}
Lemma~\ref{lem:neck} excludes overhangs whose necks live at scales
$\ge r_*$. It says \emph{nothing} about overhangs entirely below scale
$r_*$: a cap of radial extent $\le\eta$ sitting on a neck of width
$<c_1r_*$ at a scale $<r_*$ is invisible to the Serrin deficit. This is the
unavoidable $\eta$-versus-$r_*$ tension noted before:
\begin{itemize}[leftmargin=2em]
\item if $\eta\ge r_*$, all annular features are above scale $r_*$, necks
are excluded, $\Om$ has no overhang \emph{modulo a sub-$r_*$ remainder},
and Theorem~\ref{thm:realize} applies with $\eps$ replaced by the
above-$r_*$ perimeter deficit;
\item if $\eta<r_*$, the entire annular layer is sub-$r_*$ and tentacle
exclusion is blind to it; only the trivial bound \eqref{eq:trivial},
$|G\setminus\Om|\lesssim\eta$, survives.
\end{itemize}
Quantitatively, the overhangs that escape control occupy lateral scale
$<r_*$ within a layer of width $\eta$, so their total volume is at most
$C\,r_*\cdot(\#\text{necks})\cdot\eta$; bounding the number of disjoint
sub-$r_*$ necks by $P(\Om)/r_*$ gives the residual pocket bound
\begin{equation}\label{eq:residual}
   \abs{G\setminus\Om}\ \le\ \frac{(P(\Om)-2\pi)^2}{4\pi}
   \ +\ C\,P(\Om)\,\eta\quad(\text{above-}r_*\text{ part}+\text{sub-}r_*\text{ residual}),
\end{equation}
in which the first term is the genuine quadratic gain and the second is the
honest residual.
\end{remark}

\section{Assembly and an honest accounting}

\begin{corollary}[Selection-free realization, corrected]\label{cor:assembly}
Let $\Om\subset\R^2$, $|\Om|=\pi$, satisfy input (A): the density
\eqref{eq:density} and inclusion \eqref{eq:annulus} with width $\eta$ and
$r_-=\rho_i\ge\tfrac12$, at scale $r_*$. Then $\Om$ has finite perimeter
and its radial hull $G$ about $z$ is a radial graph of amplitude $\le\eta$,
with
\[
   \abs{G\setminus\Om}\ \le\
   \begin{cases}
     \dfrac{(P(\Om)-2\pi)^2}{4\pi}+C\,P(\Om)\,\eta, & \eta\ge r_*\ (\text{no overhang mod sub-}r_*),\\[2ex]
     2\pi\rho_e\,\eta, & \text{always (trivial)},
   \end{cases}
\]
and in either case $\dT(G)\le 2(\dT(\Om)+|G\setminus\Om|)$. No selection,
no $C^{1,\gamma}$, and no atlas-norm smallness is used; only the
no-overhang information extracted from $\eta$ versus $r_*$.
\end{corollary}

\begin{proof}
Finite perimeter and the inclusion come from (A). When $\eta\ge r_*$,
Lemma~\ref{lem:neck} and Remark~\ref{rem:residual} give no-overhang modulo
a sub-$r_*$ remainder, and Theorem~\ref{thm:realize} applied to the
no-overhang part yields the first bound \eqref{eq:residual}; the trivial
bound \eqref{eq:trivial} always holds. The torsion estimate is the planar
fill--rescale lemma (monotonicity of torsion under inclusion plus the
disk-torsion formula), unchanged from Plan~4.
\end{proof}

Feeding $G$ into the radial-graph theorem $(G)$ then yields
$\AF(\Om)^2\lesssim\dT(\Om)+|G\setminus\Om|+\eta^2$; the realization
contributes $|G\setminus\Om|$ at whichever of the two rates above is in
force.

\section{What is sharp, what is not, and the remaining gap}

\paragraph{Realization is quadratic only under no-overhang.} The earlier
boxed claim --- ``realization is sharp; sharpness of the whole route =
sharpness of the annulus width $\eta$'' --- was wrong, because it rested on
the false unconditional decomposition \eqref{eq:wrong}. The corrected
statement is:
\[
   \boxed{\;\text{realization is quadratic } \Longleftrightarrow \text{ no overhang};
   \quad\text{finite perimeter alone gives only the linear } O(\eta).\;}
\]
The quadratic squaring itself (isoperimetry per pocket plus
super-additivity) remains regularity-free; what it now demands as input is
the no-overhang condition, not mere finite perimeter.

\paragraph{The rate now has two carriers.} The final exponent is governed
by \emph{both} the annulus width $\eta=C(\eps^\theta+r_*)$ \emph{and} the
overhang scale $r_*$:
\begin{itemize}[leftmargin=2em]
\item if $\eta\ge r_*$, the quadratic term $(P(\Om)-2\pi)^2$ dominates and
realization gains a power, but a linear residual $C P(\Om)\eta$ remains
from sub-$r_*$ overhangs;
\item if $\eta<r_*$, no quadratic gain is available and
$|G\setminus\Om|\lesssim\eta$ is the best one can assert.
\end{itemize}
At the boundary-layer scales $\eps\sim\dT^{1/4}$,
$\delta_{\mathrm{iso}}\sim\dT^{1/2}$ one has
$r_*\sim\eps^2+\delta_{\mathrm{iso}}^{1/4}\sim\dT^{1/8}$ and
$\eta\sim\dT^{1/8}$, i.e.\ $\eta\approx r_*$: precisely the borderline,
where the residual $CP(\Om)\eta\sim\dT^{1/8}$ and the would-be quadratic
term $(P-2\pi)^2\sim\dT$ are not separated favourably. So at the rate that
(A) currently delivers, realization gives no improvement over the trivial
$O(\eta)$. An improvement requires (A) to deliver $\eta\gg r_*$, i.e.\ an
annulus width strictly larger than the density scale --- the opposite of
what one might expect, and the real obstruction.

\paragraph{The remaining gap, restated honestly.} The programme trisects:
\begin{itemize}[leftmargin=2em]
\item \textbf{Quadratic squaring:} \emph{done}, regularity-free, given
no-overhang (Theorem~\ref{thm:realize}).
\item \textbf{No-overhang:} reduces to excluding sub-$r_*$ overhangs;
tentacle exclusion handles necks at scales $\ge r_*$
(Lemma~\ref{lem:neck}), and the residual \eqref{eq:residual} is honest and
generally nonzero.
\item \textbf{Inclusion (A):} the Magnanini--Poggesi oscillation bound
under density, delivering the annulus width $\eta$ and the scale $r_*$;
crucially one needs $\eta\gg r_*$ for the quadratic realization to bite.
\end{itemize}

\paragraph{Why this matters for the original plan.} The original
obstruction was phrased as ``realization needs $C^{1,\gamma}$, which needs
selection.'' This note shows the truth is intermediate: realization does
\emph{not} need $C^{1,\gamma}$, but it does need no-overhang, which finite
perimeter does not give (Proposition~\ref{prop:lollipop}) and which density
gives only above scale $r_*$. The selection step bought $C^{1,\gamma}$,
which trivially implies no-overhang at all scales; the selection-free route
must instead win no-overhang from the deficits, and the lollipop pinpoints
exactly the configurations it must rule out. That is a sharper and more
truthful localisation of the difficulty than the earlier draft, even though
it is less triumphant: the realization half is regularity-free, but it is
not deficit-free.

\begin{thebibliography}{9}

\bibitem{AFP} L.~Ambrosio, N.~Fusco, and D.~Pallara, \emph{Functions of
Bounded Variation and Free Discontinuity Problems}, Oxford Univ.\ Press,
2000.

\bibitem{MP} R.~Magnanini and G.~Poggesi, \emph{On the stability for
Alexandrov's Soap Bubble theorem}, J.~Anal.\ Math.\ \textbf{139} (2019),
179--205.

\bibitem{FMP} N.~Fusco, F.~Maggi, and A.~Pratelli, \emph{The sharp
quantitative isoperimetric inequality}, Ann.\ of Math.\ \textbf{168}
(2008), 941--980.

\bibitem{BDV} L.~Brasco, G.~De~Philippis, and B.~Velichkov,
\emph{Faber--Krahn inequalities in sharp quantitative form}, Duke
Math.~J.\ \textbf{164} (2015), 1777--1831.

\end{thebibliography}

\end{document}
